\chapter{Introduction}
\label{sec:introduction}

Vegetation surveys require expert botanists to identify every species within a standardised quadrat, a process that is accurate but slow and limited by the supply of trained taxonomists~\cite{lifeclef2026}. PlantCLEF 2026~\cite{plantclef2026, lifeclef2026} addresses the automation of this task. Given a high resolution photograph of a quadrat, a system predicts the set of species present from a taxonomy of roughly 7,800 candidates, scored by a macro averaged per quadrat F1 that is averaged within and then across transects so that densely sampled sites do not dominate.

The benchmark is difficult for two reasons. The training images show single specimens with one label each, whereas the test images are dense vegetation scenes in which species overlap and occlude one another, so a model trained on isolated specimens transfers poorly to them. The label distribution is also heavily long tailed, which biases standard cross entropy training toward common species and depresses the macro F1.

This report describes the system submitted by the Australian National University team (ARM Wision). We partially fine tune BioCLIP 2.5~\cite{stevens2024bioclip}, a contrastive vision language model whose pretraining on TreeOfLife-10M aligns its feature space with the Linnaean taxonomy, unfreezing only the last four transformer blocks of the ViT-H/14 backbone to preserve this prior and adding auxiliary genus and family heads as regularisers. We train on an enlarged manifest of 2.65 million labelled images. At inference each quadrat is split into a $4{\times}4$ grid of sixteen tiles, encoded at $224$ and $336$ pixels and averaged, and class prior logit adjustment with an adaptive probability threshold selects the output set. The system reaches a public Macro F1 of 0.41826.
